\section{地方的尺度：不是所有边地都在等候朝廷}

\subsection{燕云的一天并不从上京开始}

燕云城镇的居民可以同时感到几种政治时间：春耕前是否要出役，哪一支驻军要补给，何日有商队经过，寺院何时募缘，家族何时为亡者立碑。上京的诏令、南京的官署和附近堡寨的需要，会在这些事情上相遇，却不一定以同样的速度传来。辽朝把幽州经营为南京之后，\eventref{LIAO-0938-EVENT-031}{南京的设置}为南部城市带来更强的中枢联系，也增加了仓储、运输、守备和文书的要求。城市不是一块被收入王朝地图后就不再变化的土地；它是不断处理这些要求的场所。

一座城的粮食不只来自城外田地。军队、官署、寺院、旅人和市场都要取用柴薪、布帛、牲畜、铁器与人力，关津和道路决定其中哪些东西能够按时进入。我们没有连续的辽南京市场账簿，不能把城市经济写成精确总量；但墓志、碑刻、钱币、城址和宋使记录使一些接口可见。它们分别来自不同的保存环境：一通墓志偏向家族的体面，一次使行偏向观察者觉得值得报告的异俗，一处遗址只覆盖被发掘的地面。把这些碎片叠在一起，不是为了制造无缝的城市肖像，而是为了看清国家命令要借谁的手才会变成日常。

辽朝的南部官署常被概括为“南面制度”，这个名称容易让人以为有一套完整的汉地行政样本。实际上，官名相同并不保证职能、人员或地方执行相同；一名官员能否征到粮、判结纠纷或组织修城，要看书吏、地方豪强、兵士与运输者是否合作。汉文文书和碑志显示汉字在这些行政与家族网络中有重要位置，却不让我们把所有居民叫作同一种人。契丹文字与其他物质线索的存在，同样不必被译成封闭的文化区。制度真正运行在多种文字、习惯和利益的交会处。

\subsection{辽河以东的山林与贡纳}

东北的河流、山林与道路让辽廷面对一种不同于燕云的治理问题。猎获、皮毛、马匹、木材与人力的流动，不能简单地按南部州县的税赋方式理解；部众首领、驻防、互市和婚姻关系可能同样重要。正史把许多关系压缩为“贡”“叛”或“讨伐”，这些词有助于排列朝廷行动，却不应替代对地方条件的解释。一次贡纳需要谁收集物品、谁保证道路、谁从中获得保护或利益，往往比称谓本身更能说明关系是否稳定。

女真诸部在辽代后期的行动，正发生在这种不均衡的网络里。后来金朝的胜利容易使读者把此前的关系看作不可避免的前奏；但对当时的人而言，辽的驻防、宋辽边境、邻近部众与本地资源都仍在变化。阿骨打可以利用已有的联系和不满组织力量，也必须把一次次军事成功转成可维持的联盟、补给和指挥。辽廷失去的不是一条抽象的东北边界，而是许多曾经使贡纳、征发和保护可以相互交换的具体接口。

考古与碑刻能帮助我们追问这些接口，却不能把材料沉默变为确定叙事。一件契丹文字题记的释读可能仍有争议，一处城址的分期也不等于当地所有活动的完整年表。最稳妥的结论不是“辽从未控制东北”或“辽完全控制东北”，而是不同地点、不同时间的控制程度不同，并且依赖地方人群是否愿意或被迫参与。

\subsection{灵州的水：一条渠如何进入国史}

灵州和河套附近的灌溉空间，是西夏早期国家化不能绕开的资源。水渠使土地能够养活城镇与军队，但水不会因为国号改变便自动流向新的仓库。修堤、开渠、清淤、分水和守护都需要劳力与地方知识；战事会截断劳力，旱涝又会改变任何计划。李氏政权能够在这一地区扩展影响，既依赖军事与外交，也依赖能否把水利、谷物、城郭和道路组织成可连续使用的网络。

宋方材料在边警和谈判时格外关注灵州、夏州等地点，常把它们作为应守、应取或应防的战略名词。对居住者而言，它们还是田地、井渠、市场和亲属所在。史料很少让我们逐户看见征发怎样落到谁身上，不能以一幅想象的农耕图填补；但也不应让水利只剩下军事地图上的蓝线。西夏的城镇可以维持，不仅因为城墙存在，也因为城外有人能在风险中继续维护使粮食到达的条件。

河套的农业和牧业也并非两个隔绝的世界。牲畜需要草场与季节性移动，农田需要水与固定劳力，军队会同时索取马匹、粮食和运输。统治者在征战、赐予、互市与修城之间分配资源时，某一选择会改变另一种生计的余地。我们无法用现存材料算出每一类负担的比例，却可以看见为什么“边防财政”必须包含水、草、道路和家庭劳力，而不只是军费数字。

\subsection{黑水城的书桌}

黑水城出土的文书让一张书桌进入西夏史。纸张上可能是佛经、契约、账目、行政手续或练习文字；这些不同文本的并存，说明书写不是只在皇帝诏书中发生。有人要制纸、磨墨、雕版、抄写、保存，也有人要在事务发生后把它变成可被追问的文字。文书中出现的官署和称谓，可以帮助我们理解制度语言，却不能让我们假定每项规定都在远方同样实行。

书桌也有缺席者。留有名字的人往往是书吏、施主、僧人、官员或能够进入文书程序的当事人；为纸张提供原料、搬运货物、在田间劳作的人则较少直接说话。多语文本不会自动弥补这种社会不平等。西夏文、汉文、藏文和回鹘文的同处，提示不同网络和读者群彼此相连，却不能证明每个人平等地使用文字。读者应把它们理解为一座城内及其周边存在多种书写需求的证据。

黑水城保存得多，正因其保存史特殊。遗址的废弃、埋藏与近代发掘或流散，使某些文本得以幸存，其他地区的档案却未必留下。研究者必须逐件核对编号、层位、语言、纪年与释文；任何从黑水城推出“西夏全社会”的结论都过大。它的力量恰恰在于局部：在这处地方，西夏不再只是宋人边报中的敌军，而是有写字的人、用法的人、供养经卷的人和处理债务的人。

\subsection{高丽方向的消息}

辽的外部关系不只由宋朝来衡量。高丽面对辽朝时既要处理边境和使节，又有自身的王统、礼仪和海陆交通；高丽材料所留下的视角，不能被辽、宋任何一方的记述吞没。不同史书对一次往来所用的称谓可能不同，礼物清单也未必完全一致。差异并不总是错误，它首先显示记录者关心的是什么：谁获得礼仪优势，谁承受军事压力，谁能控制通往都城的道路。

高丽与辽之间的使节和物资流动还提醒我们，辽东与海上并非静止的边缘。港口、河口、陆路和季节会决定消息抵达的速度；一次成功来使不能证明航线全年安全，一次冲突也不能证明所有往来停止。外部材料最适合帮助我们拆开中原正史的单一视角，却不该被当作天然中立的旁证。每部文本都产生于特定宫廷、宗教或外交环境中。

\subsection{从庆历到元丰：西北不是一个宋朝问题}

宋夏和议之后，沿边并没有进入无事可记的间歇。堡寨需要维护，互市要受检查，地方官仍会把军粮、运输和人力报到中枢；西夏也要在河套、河西和对辽关系中安排资源。一〇八〇年代的冲突和永乐城的陷落显示，较长时间的相持并不会消除前线脆弱性。每当宋廷推进据点或改变策略，西夏方面都不只是在“反应”，而是在判断如何保护水源、道路和部众的行动空间。

宋朝的政论常把西北战争同改革、财政和将领责任联系起来，这些讨论能帮助读者理解宋方焦虑，却不能替西夏写出内部动机。西夏文书与河西考古能够提出另一组问题：行政如何处理土地、宗教和书写，城镇和商路如何维系；但它们也不提供完整的作战记录。两边材料的不可互换，正是多政权历史的基本条件。战争日期可以互校，居民经验和制度实际覆盖范围则需要更谨慎的局部证据。

\subsection{天祚与仁孝：两种延续}

辽末天祚时期与西夏仁孝时期在时间上相近，却不应被写成一方衰亡、一方静止的对照。辽的东北关系、燕云防线和宫廷资源不断承受压力；西夏则在金、宋与河西网络之间维持自身的位置。两国都要处理继承、军队、宗教和地方财政，只是材料保存和政治环境不同。比较的目的不是给它们排出“强弱”，而是看到不同制度怎样面对道路、城市与多语人群。

辽灭亡后，契丹贵族、城镇居民和宗教网络以不同方式进入金及更远的世界；西夏在蒙古压力下也经历类似的分层终结。王朝终局是重要日期，但不是理解社会变化的唯一钟表。写到这些结局时，应继续询问哪一个地区被接收、哪一种文字还被使用、哪一条水渠仍被维护，避免让征服者的胜利叙事抹去地方时间。
\subsection{一座城的粮道}

辽南京、西夏兴庆府和河西城镇看起来相距很远，却有同一个难题：城内的权力如何持续得到城外的供给。粮食不是“国力”这个词自动带来的东西。它要被种出、收集、储存、记账、装车或装船，再穿过会受季节、战争与检查影响的道路。仓门关上时，城中的官员可以讨论边防；仓门打开时，承担搬运的人、管理牲畜的人、守关的人和等待配给的人才让边防有了具体形状。

辽的多都城安排使供给问题格外复杂。上京的政治仪式、南京的南部资源、东北和草原的军政关系，需要由不同方向的道路和人员维系。辽史中的中央调度能够提示这种规模，却很少让读者知道某次运输在何处延误、谁因此改变耕作或贸易。城址和交通研究不能补出这些细节，却可提醒我们城墙、河谷和道路坡度都是行政要面对的物质条件。一次军队的到达，不表示下一次粮食也能到达。

西夏的粮道同样要经过水。兴庆附近的灌溉、河西绿洲的田地和通向边寨的道路，彼此不能替代。战时，官府可能要求更多粮食、人力和牲畜；平时，修渠与维护城堡也会占用劳力。史料没有给出完整的征发账，因此不能把任何一项中央数字换算成农户的真实负担。然而，正是这些缺口让我们看见问题所在：一个国家若把城堡推进到更远处，便把随后的运输、取水和维护责任推给更长的一串地方。

\subsection{经书的旅行与工匠的手}

一部经书从抄写或雕版到被放进寺院，并不是抽象的“文化传播”。它需要纸、墨、木材、刻工、校对者、施主和能够保存它的建筑。辽代佛教材料、西夏的印本与文书，说明不同政权下都有把宗教文本变成物件的实践；它们也提示工匠和僧侣必须在城市、寺院、市场与道路之间取得材料。某些文本使用汉文，某些使用西夏文或其他语言，不应被理解为彼此封闭的世界，而应追问读者、赞助者和使用场合。

宗教物品可以跨越边界，但其旅行从不无条件。和平时的使节和商路可以提供接触，战争时的检查、劫掠和封锁会改变风险。即便同一政权内部，远处寺院是否能获得纸墨，也取决于地方市场、捐施和保护关系。艺术史或文献目录常把作品按风格、语言或朝代分类；放回行动中看，分类之间还隔着运输、翻译、抄写和资金。那些环节很少在正史中被连续记录，却使多语宗教世界成为可能。

寺院也并非只有施与的一方。它可能接待旅人，组织仪式，保存文本，拥有或管理某些资源，并与地方家族和官府建立关系。保存下来的碑刻偏向有能力留下名字的施主，不能据此判定普通居民是否分享同一种信仰或经济收益。辽、西夏的宗教史不需要被写成国家“利用”宗教或民众“信仰”宗教的二选一；较好的问题是，在特定城镇和年代，谁用哪些资源让一座寺院、一部经书或一次仪式继续存在。

\subsection{名称之外的身份}

宋、辽、西夏、金的史书不断使用契丹、党项、汉人、渤海、女真、蕃部等名称。名称是政治沟通、行政分类和叙事安排的重要工具，但它们不会自动给每一个人规定稳定而唯一的身份。一位家族成员可以以祖籍、官职、婚姻、居处或所用文字进入不同文本；同一个地名也可能在不同政权的记录里有不同写法。将这些复杂关系压缩成现代式、彼此排斥的标签，会使迁居、服务、通婚、贸易和宗教往来都失去位置。

这并不意味着差异不存在或权力没有边界。官署的分类可以决定谁服何役、哪类人可进入何种程序；战争和征服会使名称带来危险或机会。问题在于，我们必须先弄清名称出现在哪一种材料、为谁服务，再讨论它的社会后果。墓志说一个人来自哪里，可能是在安置家族荣誉；官府文书的称呼，可能是为了税役或司法；宋使所写的异俗，则反映使者的观察位置。每一种证词都能说明一点，不能独自说明全部。

\subsection{两种终局后的地方时间}

辽在一一二五年、 西夏在一二二七年的王朝终局是不可替代的年序节点。它们分别改变了皇室、官署和军事资源的最高层配置，却没有让所有地方以同样速度进入新制度。辽的城镇、家族与寺院要在金的统治下重新寻找保护和文书接口；西夏的河西道路、水利和宗教网络则在蒙古统治及其后续安排中经历不同变化。征服者的史书喜欢把结局写得迅速，因为它要确认新王朝的合法起点；地方材料提醒我们，转换常常拖得更久。

这不是要把灭亡说成没有意义。战争可以毁坏城镇、夺走劳力、打断运输，并让文书和宗教场所失去旧有的保护；幸存者的选择也会受新政权的征发、分类和军事部署限制。只是“辽亡”“夏亡”不能取代后续的问题：谁能回到田地，谁保住了寺院，哪种文字仍被使用，哪些道路被重新打开。答案在不同地点不同，且许多已无法由现存材料恢复。保留这种不均衡，正是让多政权边地不再被胜者年表吞没的方式。

\subsection{史书的桌面与遗址的地面}

《辽史》的本纪、志书与列传把辽朝留在一张由皇帝、官号、战争和制度组成的桌面上；《宋史》、宋人编年和使行录又把同一北方写成宋朝必须处理的外部。它们对于日期、使节、改元和军政决策不可或缺，却都会把地方活动压缩为朝廷注意到的部分。读者不必因此放弃正史，而应在使用它时知道它的焦距：它善于说明中枢何时作出何种处置，不善于说明处置在每一条道路和每一户家庭如何发生。

城址、墓志、寺院遗存和黑水城文书把镜头移到地面。城墙的修筑、墓主的官历、题记中的施主、纸卷上的手续，都有更具体的地点和物质痕迹；但地面的材料也有自己的盲区。发掘区不等于整座城市，墓志首先是精英记忆，遗址的废弃层与器物的使用时间未必相同，文书的保存很可能来自偶然。把它们和正史并读，并非让一种材料“纠正”另一种，而是让不同的证据回答不同的提问。

例如，辽南京的经营可以由帝纪和官制材料确定大致节点，城市遗址和碑志则让我们追问谁在这座城中拥有可见的资源；西夏国家的年号和对外战争可以由宋辽金元史互相校正，黑水城的文书、印本与寺院材料却使官府和信仰在局部实践中显现。两种材料都不能直接算出社会全貌。历史叙事的诚实，不是把这种差距藏在注释里，而是让读者看见：我们为何能确定某个日期，却不能确定某位无名运输者在同一天做了什么。

\subsection{读者回到道路上}

从潢河到幽州，从灵州到黑水城，辽与西夏的历史并不沿着一条朝代更替的直线展开。它由不同的人在不同地点一次次选择：统治者选择是结盟、出兵还是遣使；官员和书吏选择用哪些文字、名目和程序处理事务；商旅与地方家庭选择何时出行、把货物卖到哪里、在战事来临时依靠谁。材料不能让我们替他们说出内心，却能让我们看见选择受到水、路、仓储、亲属、寺院与政治名义的约束。

因此，下一次在史书中读到“盟”“贡”“破城”或“亡国”，不妨先停在一座关口、一段渠或一张文书前。谁能通过，谁要交纳，谁能读懂，谁必须搬运，谁在终局之后还留在当地？这些问题不会把大事件缩小，反而使辽和西夏不再只是宋朝叙事两侧的名字。它们是拥有自己的时间、语言和地方社会的行动者，而这些社会的残缺证据仍要求我们谨慎而持续地阅读。
\section{在边界上办事的人}

\subsection{馆伴的行囊}

辽宋使节的相遇，通常由盟书、国书和皇帝的称谓开始；真正让相遇发生的，却是一支必须带着衣物、食物、印信、翻译和随从穿越关津的队伍。馆伴要知道哪一处渡口可用，哪位地方官负责供应，哪一种礼仪会在抵达时被要求。他的差事不是单纯传话：一项称谓写错、一次马匹供应延误、一次随从越过规定区域，都可能被中枢解释为对盟约的试探。澶渊以后，\eventref{LIAO-1005-EVENT-038}{盟约的落实}正是在这些看起来琐碎的事务里被反复确认。

宋人使辽行录保留了许多馆舍、饮食、服饰和仪注，因而读来像一扇可以打开辽社会的窗。窗的框架却由使者自己的任务决定。他尤其注意座次、仪仗和让自己感到陌生的语言，却未必记录替他添草料、修车轮和传递文书的人。辽方材料更少留下同一行程的镜像。把两类沉默并置后，读者可以把使团当作一项基础设施来理解：它需要国家授权，也需要沿途城镇愿意承担日常成本；它可显示政权关系的温度，不能替代普通边民的完整声音。

西夏使节面对的道路更不稳定。兴庆府、陕西沿边、河西与辽境之间的通道，会因战事、封市、洪水或地方冲突而改变。使者携带的礼物和文书不仅是王权的象征，也是让接待者能够把这支队伍识别为“可交涉的人”的凭证。宋方记录中的请封、朝贡或谢恩，常把西夏使节置于本朝礼制；西夏的年号、官署与文字则表明它并不只在宋朝的词典里行动。我们很难从现存材料复原一趟行程的每个夜宿点，但能看见外交必须依赖一批熟悉地名、语言和危险的人。

\subsection{关市的早晨}

开市前，关市并不是一片自然形成的热闹。守门者要确认日子和禁令，商人要判断带去的是布帛、茶、盐、皮毛还是马匹，地方官要处理税、检查和纠纷。辽宋边市的和平价值不在于取消边界，而在于把相遇压进可以预期的时间和地点。马匹、铁器和消息的流动仍会引起军事担忧；一处市场同时是收入来源、情报节点和潜在风险。对携货的人而言，一道临时禁令可能比远方朝廷的盟约更直接地改变一年生计。

西夏与宋的互市也不应被写成贡赐制度的附属物。河套、陇右和河西的居民在城寨、牧地和绿洲之间交换物品，官方允许、禁止或征税的对象，会影响哪些交换能留下文书。史书往往在战争、禁市和和议时才让市场进入叙事，平常日子里的小额交易则几乎无声。钱币、窖藏和遗址能补出一些物的移动，却不能替任何一位商人说明价格、利润或风险。市场史的困难正是它的价值：它迫使我们把国号与实际生活之间的关津、度量和信用写出来。

市场也带来语言问题。讨价还价、核验货物、解释禁令和解决债务，都需要有人能在不同语词之间使意思暂时对得上。翻译者可能是官员、商人、僧侣、地方首领或同时承担几种身份的人。材料很少让他们以自己的口吻出现，然而没有他们，条约中的“互市”无法变成可执行的安排。将多语世界只写成几种文字并列，会错过这些口头工作；将它写成毫无摩擦的交流，又会忘记误译、欺诈和权力不均也会沿同一条道路移动。

\subsection{北面与南面之间的账簿}

辽朝的北面、南面官制常被用来说明一个草原政权如何统治农耕地区。这个解释若停在名称上，反而会把办事的人抹去。官署需要能够登记、传达、裁判和催纳的书吏；军队需要知道粮草何时到达；地方豪强、寺院和部众首领会用自己掌握的资源影响命令的落点。北面与南面不是两只内部完全一致的盒子，而是不同地区、不同来源的资源被朝廷尝试接合的方式。

一份制度条文可以告诉我们官名和职掌，却不能证明每一处官署都按同样方法工作。辽的百官志、营卫志和食货志本身也是后出的整理，它们把长期变化压缩为较稳定的制度图像。契丹文字题记、汉文墓志和城址材料能将某些官名、家族和地点重新钉在地面上；它们又偏向能刻石、能安葬、能被发掘的少数人。读者不应在“中央文献”与“地方材料”之间选择一方为真，而应让二者互相限制：前者说明制度想要什么，后者偶尔说明谁能够把它变成事实。

西夏的官文书和法律材料提出了相似的问题。法卷可以规定如何分类财物、罪名和程序，实际的案件却需要当事人、书写者和裁断者都进入同一套手续。黑水城保存的材料使这种书写世界可见，却不能让我们推断全境都有同样密集的官署与识字者。西夏本国连续国史未存，王朝政治的年代常要由宋、辽、金、元史互证；正因如此，一份出土文书的局部性不是缺陷，而是防止敌国叙事独占解释权的必要边界。

\subsection{一座寺院如何穿过战争}

寺院的屋顶、经卷和造像看上去远离边防，实际上也要面对道路和资源。修一座塔需要砖木与工匠，抄一部经需要纸、墨和能读写的人，维持僧众则需要食物、土地或捐施。辽代的寺院题记、经幢和墓葬图像显示佛教如何同家族记忆、城市声望和政治保护相连；它们不证明全体居民共享同一种信仰，也不说明国家能无条件支配寺院。宗教网络之所以值得进入编年，是因为它提供了战争、迁徙和政权更替之外仍在运作的资源与文字通道。

西夏佛经印本尤其提醒人们，宗教不是只有口头仪式的领域。翻译、雕版、抄写和供养要把多种语言、技术与资金聚在一起。汉文、藏文、西夏文和回鹘文文本在同一地区出现，不等于每个居民都能读写每一种语言；它更明确地告诉我们，寺院、官署与商路需要处理不同读者和不同权威。某一部经书的保存，说明有人曾使它成为值得保存的物，不说明它代表西夏所有地区或所有宗派。

战争会让寺院的保护关系改变，却未必让宗教生活立即停止。辽亡后，原有寺院可能在金的官府、地方家族或新的施主网络中寻找继续存在的方法；西夏在蒙古征服后也有文本、图像和人员以不同方式进入新的政治环境。征服史书关心城破和君主投降，往往少写一处寺院何时重修、一批经卷怎样转移。不能由这种沉默虚构连续性，但也不能把王朝终局写成所有书写与祈愿同时消失。

\subsection{墓地里的迁徙}

辽代汉文墓志和契丹文字材料常为后人安排一个人的祖先、官历、婚姻与葬地。它们有意把家族写成有秩序的延续，追赠的官衔和理想化的品行不能直接当作生前事实；然而，正因为要解释一个人为何属于某个家族，迁徙和任职的路线也会露出来。祖籍、任所和埋葬地不在一处，提示道路、军役、婚姻和官署能把家庭拉向不同的中心。

西夏墓葬和题记的保存更不均衡，不能拼出同样连续的家族地图。黑水城文书中偶见的姓名、债务、施主或官署关系，能够把抽象的“党项社会”拆成若干具体当事人，却很少告诉我们其一生如何度过。女性、儿童、雇佣者、牧人和流动者尤其难以在史料中持续追踪。承认这些人难以被听见，不是把他们逐出叙事；相反，国家征役、修渠、战争运输和市场禁令之所以重要，正因为它们会落到这些没有留下墓志的人身上。

\subsection{女真与青唐：第三方改变所有人的选择}

辽宋关系在一一一四年之前已不是两方独自安排的边界。女真诸部的行动、东北的贡纳与贸易、宋廷获得的消息，都让燕云和河北的计算发生变化。宋人把女真来使写成边政消息，辽廷则要处理原有驻防和部众关系的松动；女真首领也在不同邻居之间寻找资源。后来的金朝兴起不应被倒灌成一种天命，它是许多道路、据点、部众和政治承诺重新组合的结果。

西夏面向的第三方同样不止宋和辽。青唐、河西绿洲、吐蕃诸部与回鹘社群构成不同的贸易和军事空间。西夏向甘凉方向的行动会触及城镇、水源和宗教网络，而不只是扩张一条地图边线。宋方史料常从本朝边防看待这种变化，西夏文书则更容易显示某一城镇的行政与宗教实践；两者之间仍有大量未被保存的地方经验。外部世界不是一张附录地图，而是促使辽、西夏统治者反复调整联盟、驻军和使节的压力来源。

\subsection{河流冻结时}

北方的道路不是全年可用的直线。河流封冻可以让某些路段更易通过，也可能让补给和渡运变得危险；春汛、泥泞、风雪和缺水会把军队、使团和商队的时间表打乱。辽军南下、宋军北伐、西夏的边地调动，都必须处理这一层季节性。环境不是决定胜败的独立主角，同一条河、同一段山路会因准备、地方合作和战争局势而产生不同结果；但任何忽略它的叙事都会把人和物写得像在空地图上移动。

灾异记录必须更谨慎地使用。正史在本纪中记录旱、蝗、洪水和地震，常与政治劝戒相连；一条记录可说明朝廷注意到某种危机，不能自动证明整个政权的人口或财政同样受灾。水利遗址、沉积和地方材料能补充环境史，也有各自的时间尺度。对辽与西夏而言，较可靠的问题不是“自然是否导致衰亡”，而是某次水、路和气候的变化怎样让既有的征发、储粮、迁徙或防守更难维持。



\section{河西的水渠、寺院与一张无法归一的国史}

西夏的都城与河西城镇并不只靠军队相连。黄河灌溉、泉水、渠堰、屯田、市场和通向绿洲的道路，决定一座城能养多少人、能留多少军马，也决定朝廷命令在何处会被季节和水源改写。宋方记载常在战争爆发时才特别注意这些地方，把它们称作边寨、险要或进军路线；对当地居民而言，水渠首先关系到一季耕作、牲畜饮水和城内粮价。把西夏写成宋朝西北的军事投影，便会漏掉这种以灌溉和交通维持的国家日常。

黑水城及河西出土的西夏文、汉文、藏文和回鹘文材料，使这种日常偶尔露出轮廓：经卷、契约、官文书、印本和题记并列保存，提示书写本身已在行政、交易和宗教之间来回流动。它们不是一部完整的西夏档案。出土地点集中、文书定年不齐，许多文本的流传过程也不明；一张账目不能替全境计算赋役，一件佛经不能代替所有人的信仰。不过，正因为这些残片不是宋廷使者的眼睛，它们迫使我们把河西看作有自己文书、寺院和市场节奏的地方。

寺院在西夏并非战争叙事之外的宁静角落。抄经、雕印、供养、修建和僧侣流动需要纸、木材、颜料、劳动和捐施，也需要相对可通行的道路。统治者借佛教仪式和译经建立权威，地方施主则可能通过造像和题记维系家族、职业或社区的名望。宗教网络能够跨越政权边界，却不等于它天然反对或支持任何朝廷。材料可以让我们看见某一寺院、某一批印本的活动，不能由此推出所有西夏居民共享同一套信仰或政治态度。

仁宗时期的西夏处在金、南宋、吐蕃诸部和河西商路交错的环境中。对外关系中的贡赐、册封与互市，是统治者调节边境的手段，不是单向度的臣属符号。马、茶、织物、药材和贵金属的流动，要经由关口、商人、翻译者与地方征收；一纸和约若不能让商队安全抵达、让城寨得到粮草，便难以在日常中成立。宋、金与西夏关于称谓的文字特别容易被后来史家读成等级秩序，实际的交换却总在具体道路和价格中协商。

西夏的军政压力也会回到水利和家庭。城寨需要守军，守军需要粮食，粮食来自有水的田地和能够把谷物送到城内的运输。加筑一座堡垒可以暂时扩大前沿观察和驻军能力，却可能使周边村落多出服役、运粮和维修的义务。官修史料常将这种变化写成“筑城”或“出师”，很少留下渠工、车夫和妇女如何调整家计。承认这一沉默，并非把边防成本化为无从讨论的背景，而是拒绝用一支想象中齐整的西夏军队遮住它。

辽末与金兴起后，西夏的选择更难从“亲谁”二字概括。辽、金、宋的力量转换，改变的是河套、河西与陕西各条道路上的使节、商队和军队能否通行；西夏朝廷必须在旧关系、新威胁和内部资源之间留出余地。外交使团或许在都城完成仪式，边地城镇却要等待货物、文书与保护何时真的抵达。正因如此，一次受封、一次联姻或一次出兵只能说明某一时刻的政治安排，不能替代对地方控制和民生后果的判断。

蒙古军进入河西前后，旧有多语世界并未因为战报中的“攻下”而立刻消失。征服者需要向导、粮食、工匠和能够读写的中介；被接收的寺院、城镇和家庭也要判断新军的征发、保护与宗教承诺能否兑现。成吉思汗去世与西夏覆亡都发生在一二二七年，日期上的并置很醒目，却不能把两件事写成一幕毫无余波的终场。国号消失后，水渠仍需维护，商路仍有危险，书写与礼仪仍在新的权力关系中继续。

西夏材料的残缺要求一种更慢的叙事。我们能确认政权更替、部分城镇和战事的顺序；不能确认所有居民何时改用何种身份，也不能把宋方或元方的胜利叙事替作河西的全体记忆。留住这层不确定，反而使国家的形状更清楚：它不是地图上均匀涂开的领土，而是由水、城、寺、道路和不断重新谈判的人际关系连接起来的区域秩序。

\section{四京之间：辽的命令怎样离开帐幕和城门}

辽的多都城不是地图上的四个圆点。上京临潢府靠近契丹贵族和北方草原的活动范围，南京幽州面对燕云农业区、市场和通向宋境的道路，东京与中京又各自连接东北、山地与内地。皇帝巡幸、官员迁调和使节往来，使这些城市在不同季节承担不同的政治重量。把辽的统治想成一条从上京直落各州的命令线，会错过它必须在营帐、宫城、仓场和关津之间反复转换的事实。

南面官、北面官的名目常被后世当成两个整齐分开的抽屉，仿佛制度一立，契丹、奚、渤海与燕云居民便各自被固定安排。志书确实保存了官号和某些沿革，却较少让人看见一个诉讼、一项征发或一次迁调在地方究竟经过谁的手。官名可以说明朝廷试图如何分类，不能证明每一地都按同一方式运作。对一个驻在南京的官员、一个沿路递送文书的吏员或一户负担运输的人家而言，真正重要的或许是下一道命令来自何处、能否找到马匹、仓中有没有粮，而非制度图上为他们预留的位置。

幽州的城门和市场让这种差别具体可见。燕云并非辽廷南方边界上一块静止的领土；农田、盐铁、手工业和通向河北的道路，使它既能提供税粮，也会在宋军北进、使节通行或边市受限时首先感到紧张。城内商人、工匠、寺院和军人不必拥有同一种政治语言，却共同依赖道路和治安。辽、宋史书在谈判或战事时才格外注意幽州，城址与墓志则偶尔保留当地官员、家族和信仰网络的片段。二者合起来仍不足以写出城市全貌，但足以拒绝把南京只当作战争地图上的一个箭头终点。

圣宗、萧太后面对宋朝北伐和河北前线时，调动的也不只是骑兵。边军需要草料和粮食，驿路需要马匹与看守，渡口和城寨需要有人传递消息。九八六年宋军多路行动，双方史书各自记录胜败和将领，兵数与伤亡却难以互相校正。较可靠的场景是，燕山南北的道路突然同时被军队、运粮者、使者与避兵者使用，原本可供市集和农时往来的路径变成风险极高的通道。辽的胜利或宋的退却，不能替每一处居民说明他们怎样保存牲畜、谷物和亲属。

澶渊之后的和平尤其能照出这种多中心的成本。使节、银绢、边市和巡逻把敌对关系转成可计算的程序，却仍要依靠河道、关口和可信的中介。宋方的岁币用语、辽方的礼仪称谓都服务于本方的政治体面；如果只从其中一边判断臣属，便会把互相需要的供给安排误写成单向关系。边市的开放也不是一扇永远敞开的门。货物能否越境，取决于地方官是否执行、道路是否安全、价格是否可接受，以及战争传闻有没有使商人停下脚步。

契丹文字、汉文墓志与佛寺题记让辽的社会空间出现另一种时间。碑刻常写功德、官衔、亲属和日期，文字材料也会因保存和释读困难留下空白；它们不能按现代族群标签把人的身份一次定死。不过，多种文字和礼仪在同一区域留下痕迹，至少说明统治并非只凭一种语言完成。一个家族可以在官署、市场和寺院之间采取不同的书写与交往方式，一名被追赠的官员也未必等于实际掌权者。史料的细部正好限制我们把辽的制度描成毫无摩擦的二元结构。

辽末的压力也应从这些接口看。女真势力在东北扩大时，朝廷面对的是对部族、城镇、道路和粮源的同时失去把握；宋金之间的联盟又让燕云的防务与财政不再只是一条南方边线。天祚帝的行动在《辽史》中有清楚的王朝年序，地方上哪些城先断粮、哪些官员还能传令、哪些家庭向何处迁徙，却远不清楚。不能拿后来的覆亡替每一次当时的选择判决，也不能让“辽亡”一语遮住不同区域先后失去保护的过程。

耶律大石西行把这种不均衡延长到更远处。西行者带走的并非一个完整辽朝，而是人、记忆、政治资源与跨草原道路上的可能性；留在燕云和东北的人则在金的接收中面对另一种制度与征发。西辽的形成需要中亚、波斯文与地方材料来并读，不能只从中原王朝的余波理解。王朝终结在原有核心区可以有明确日期，制度、语言、贸易与亲属网络的转换却沿着不同道路发生，快慢不一。

\section{兴庆府外：西夏的国家在水、纸与马之间}

兴庆府附近的灌溉并不只是西夏富庶的背景。黄河与渠堰决定田地何时能播种、城中谷物从何处来、军队在旱年还能停留多久。修渠、清淤、分水和防洪需要劳力、地方协调与持续的维修；一处渠口被毁，影响的不只是农户收成，也会改变税粮、马料和城内价格。宋方史料在战争时常把这一带看成进军目标，西夏文文书和考古则偶尔让行政、宗教和日常书写出现。两种材料的视角不同，恰好提醒我们不要把水利直接等同于某一方军略的附属。

李元昊称帝前后，礼制和服饰的改变在宋方叙述中格外醒目，因为它直接挑战宋廷对名分的安排。可是对兴庆、夏州和河西的居民，国号和冠服能否成为日常政治，取决于赋役、军队、贸易和地方首领的关系是否随之调整。称帝使宋、辽与党项之间的外交词汇改变，也促使西夏朝廷更明确地组织王权象征；它不能让我们断言每个部众、城镇或绿洲已在同一天接受同一种认同。国家形成不是礼仪一刻完成的事情，而是把水源、道路、文书和人力接入更稳定秩序的漫长劳动。

三川口、好水川和定川寨的战事，必须放在高原、河谷与堡寨的尺度上理解。守军需要从后方运粮，进军者需要找到水源和可通行的道路，山地气候会令同一条路线在不同季节具有完全不同的代价。宋方兵书和战报可能特别强调将领的失误或敌军的奇袭，西夏方面没有保存与之对称的连续军令；因此我们不能把任何一方的军数、行军速度和战果变成无误的统计。可以确认的是，战争反复迫使边地社会承担筑寨、运输、征发和避兵的负担。

一〇四四年的议和没有让这种负担消失，而是把部分冲突改写为贡赐、互市和边境管理。和议为宋朝和西夏都留出调整军费的空间，也让沿线城镇有可能恢复部分交换；但商队、牧民和驻军面对的安全并不会因条约文字而自动一致。边市若要运作，需要关津、翻译、度量和地方官能够处理纠纷；一旦战事传闻再起，同样的道路又会被封锁。和平在这里并非没有军事，而是用驻防、巡逻与交涉维持一个脆弱的可通行性。

西夏的法律、佛经和印本使我们看到另一种国家能力。文字可以把土地、财物、职务和仪式记入可转交的形式，寺院也能通过供养、抄经与造像聚集资源与人际关系。可是保存下来的文书主要来自特定地点和特定收藏史，不能视为完整档案。它们足以说明河西并不只是宋军笔下的荒边，却不足以让研究者按文书数量重建全国财政，更不能从一座寺院的题记推出所有人的宗教实践。材料的缺口是叙事必须承认的地形。

仁宗时期与金、南宋及吐蕃诸部的交往，使西夏的选择不只在宋夏边界上发生。马匹、茶、织物、药材和经卷沿不同道路移动，地方中介同时处理市场、贡赐和安全。一次使节往来可以显示双方愿意暂时谈判，不能证明道路上的所有人都从中受益；一纸册封也不能将多种政治关系折成唯一的上下级。西夏保持行动空间的代价，正是不断在更强邻国与内部资源之间调整，而这种调整会落到城镇的征收、商旅的通行和边军的粮草上。

蒙古进攻前的河西更不应被写成等待征服的空地。绿洲城市、寺院、书写传统和商路仍在组织生活，蒙古军队若要取得粮食、向导、工匠和税源，必须与这些既有关系相遇。一二二六至一二二七年的战争及其破坏可由多方材料确证基本顺序，至于具体城镇的损失、迁徙规模和民众反应却没有完整记录。国号终结之后，水渠仍需要修理，僧侣和商人仍要决定依附或离开；新统治者继承的不是一片空白，而是一组需要重新维系、且已经被战争损耗的接口。

\section{边市的早晨与边报的傍晚}

澶渊议和之后，边界并没有成为两种世界之间无人经过的线。河北和燕云之间仍有道路、河流、关口与市场，商人、使节、逃亡者和巡逻军队在不同时间使用它们。早晨打开的关市要核对货物、度量和人员，傍晚传来的边报却可能让同一条道路突然不能通行。地方官既要维持交易带来的税收和供给，也要防备消息、马匹和武器越过边界。和平把战争的危险压低，却把管理可通行性的责任留给更靠近道路的人。

宋人使辽的行记常注意服饰、宴饮、制度和沿途风景，这些观察珍贵，却首先是使者在外交任务中的视线。它们能让我们知道道路经过哪些城镇、接待如何安排，也会以“异俗”给复杂社会贴上方便的标签。使团见到的是被安排给他们看的驿馆、官员和仪式，不是所有牧民、农户、商贩和军户的生活。将行记同墓志、城址、契丹文字和高丽材料交读，不能消除偏见，却能阻止我们把一位宋使的惊异当作辽境社会的全景。

辽朝向宋方输送或接受的银绢，也不应只被理解为抽象的国库数字。织造、收集、运输、验收和再次分配需要许多环节；任何一段河道结冰、道路受阻或地方征收失当，都可能改变物资真正抵达的时间。辽廷能够把这些礼物转入军政和宫廷体系，宋廷也用它们换取边境较低的战事风险，但两边都无法仅凭一纸盟约免除运输和监督。数字在史书中显得整齐，负担却在不同地方分散发生。

边境居民面对的选择尤其难以用国家立场替代。有人靠互市出售牲畜、布帛或日用品，有人因驻军获得短期生计，也有人必须承担修路、递送或搬运。战事传闻出现时，能够带走货物的人与只能守着田地的人所受的风险不同。现存材料很少让这些人直接发言；不能因此为他们虚构同一种恐惧或忠诚。比较可靠的是把他们置于可见的空间条件中：关门何时关闭，渡口能否使用，军队向哪一处索粮，都会比都城的礼仪更快改变一户人家的日常。

西夏与宋、辽的互市也有相似而不相同的难题。河套和陕西边地的马匹、茶叶、盐、药材和织物，沿着城寨与河谷移动；山地道路、旱情和军队征调会使交易突然昂贵或中断。朝廷的禁令与开放命令说明政治选择，却不能说明每一队商旅是否真的抵达。西夏有时需要通过交换获得难以自产的物资，宋方也要借市场减少边境冲突的成本；这种相互依赖从不意味着两边已经放弃防备。

边市由此是一种小而持续的外交。它并不取消战争，也不只是国家仁慈施予的贸易机会；它要求地方官、译者、商人和驻军每天确认规则还是否有效。一旦一方加强堡寨、另一方传来出兵消息，先改变的往往是货物、马匹和人的流动，而不是史书中的正式关系。把这种脆弱的日常放进辽夏的年序，才能理解为什么一个长期并立的秩序既能维持数十年，又始终在道路和仓场的细部里积累摩擦。

\section{辽末的东北：森林、河流与归附的选择}

辽与女真诸部的关系长期通过贸易、封授、军事合作和地方首领的竞争维系。东北并不是一张等待中枢命令填色的边地：河流、森林、猎场、城镇和通向高丽的道路，使不同群体有不同的资源与消息来源。辽廷能够授予官号、征调军队或要求贡赋，却难以在每一个季节同样掌握这些关系。女真力量的增长因此不能只以“反叛”解释，也要看到旧有的贸易和军政接口在何处失去可信度。

阿骨打起兵后，辽朝面对的不只是一个新军事集团，还要判断东北各城、各部和各条道路会怎样回应。官修史书常把归附写成首领的单一决定，实际的变化可能经由亲属、市场、战败消息和补给困难缓慢传开。一个城镇是否转向，不仅关乎政治名义，也关乎粮食、皮货、马匹和人员能否继续流动。材料不足以复原每一处的谈判，却不允许把复杂的区域关系写成无声而瞬间的倒戈。

辽廷调兵应对时，草原与城市的供给差异变得突出。骑兵需要草场和马料，城防需要粮仓、工匠和可通行的道路，远方命令则依赖驿站和地方官。任何一环受阻，都可能使看似充足的兵力无法在需要的地点出现。东北战事的细节在辽、金和宋方叙事中多有互相矛盾处，尤其不应照录战功和人数；但城镇、河道和牧地被争夺的事实足以说明，辽末的危机是一种基础设施与忠诚同时被拉扯的危机。

高丽方向的材料为这种区域性提供旁证。使节、贸易和边境消息从来不只服务于一个王朝的年表，观察者也以自己的政治位置选择记录什么。外部史料不是天然中立，却能使中原与辽金正史之外的联系显现。把它们同契丹文字、城址和汉文墓志并读，不能消除年代和地名的难题，但能避免把东北只写成宋辽金三方外交的背景布。

辽亡之后，东北的变化仍没有停在一一二五年。金朝接收的是城镇、人口、道路与旧官署，也接收了因战争而断裂的家庭、市场和地方记忆。有人进入新的军政体系，有人向其他地区迁移，有人继续依赖熟悉的河道和交换网络。后来的金代制度不能让辽代的社会经验自动消失；同样，辽代的官号和族名也不能完全决定新时期的身份。王朝更替在地方总是以不均匀的速度发生。

\section{河西的文本如何越过战争}

黑水城保存的文书、印本和佛教材料使西夏后期不再只能由敌手的战报叙述。它们让读者看见书写如何进入法令、经卷、账目和日常事务，也让人看见保存的偶然：一处遗址留下大量纸张，并不表示其他地区没有自己的文书，只是没有同样的存留条件。研究者可以用这些材料提问行政、宗教和语言接触，却不能以出土密度替整个国家计算社会结构。

印刷和抄经需要纸张、木板、劳力、识字者和可以持续供养的机构。战争与政权更替会威胁这些条件，也可能使某些文本、技艺和书手随人迁移。西夏的佛教世界因此不能只被理解为王室装饰；寺院、施主、工匠和旅行者构成的网络，既能维持地方记忆，也会受到税役、军队和道路安全的约束。题记记录的是愿望与功德，不是全体信众的意见，却把这种网络的物质条件保留下来。

蒙古接收河西时，面对的正是这些仍在工作或已经受损的关系。新军需要粮食和向导，行政需要会写文书的人，宗教机构则要重新与保护者协商。征服可以改换城门上的旗号，不能在一夜间替代所有渠工、书手和商人的知识。这个接续过程的材料尤为不完整，故不能把河西写成快速、完全的整合；它更像在战争留下的损耗中，反复把水、路、文字和人重新接到新的征收与保护关系。

\section{一座边城的账房}

边城的账房把远方的战争变成近处的谷物、马料和人名。守军人数、仓粮出入、修城工料和使节给食都要写进可以上报的格式，可是账房记下的数目从来不等于资源已经安稳。谷物可能还在路上，马匹可能因草场或疫病减少，承担运输的人也可能先要照料自己的田地。宋、辽、西夏的官府都试图让边防变得可计算，边防的现实却不断用季节、道路和地方关系打断这种计算。

对城外村落而言，军需并非只有一项税目。一次运粮会占去车马和劳力，修寨会改变取木、取水和耕作的时间，驻军增加又会推高市场对粮草的需求。有人可以借与军队交易获得收入，有人则因同样的征调失去余力。史料很难给出完整的受益与受损名单，不能由此把边地社会写成只有受害者或只有机会者；应当看见同一套防务如何在不同家庭和不同季节产生不相同的后果。

西夏与辽的长时段并立之所以可能，部分正在于这种日常的协调反复奏效：城寨还能取得粮，驿路还能送出消息，市场还能在有限范围内开放。它的脆弱也在这里。只要水源失常、道路被截或中枢政治改变，原本分散的负担就会同时压向仓场和居民。大事件的转折因此不是凭空发生，它常先表现为这些细小接口再也难以同时运转。
